% Options for packages loaded elsewhere
\PassOptionsToPackage{unicode}{hyperref}
\PassOptionsToPackage{hyphens}{url}
\PassOptionsToPackage{dvipsnames,svgnames,x11names}{xcolor}
%
\documentclass[
  10pt,
]{article}
\usepackage{amsmath,amssymb}
\usepackage{iftex}
\ifPDFTeX
  \usepackage[T1]{fontenc}
  \usepackage[utf8]{inputenc}
  \usepackage{textcomp} % provide euro and other symbols
\else % if luatex or xetex
  \usepackage{unicode-math} % this also loads fontspec
  \defaultfontfeatures{Scale=MatchLowercase}
  \defaultfontfeatures[\rmfamily]{Ligatures=TeX,Scale=1}
\fi
\usepackage{lmodern}
\ifPDFTeX\else
  % xetex/luatex font selection
  \setmainfont[]{DejaVu Sans}
  \setsansfont[]{DejaVu Sans}
  \setmonofont[]{DejaVu Sans Mono}
\fi
% Use upquote if available, for straight quotes in verbatim environments
\IfFileExists{upquote.sty}{\usepackage{upquote}}{}
\IfFileExists{microtype.sty}{% use microtype if available
  \usepackage[]{microtype}
  \UseMicrotypeSet[protrusion]{basicmath} % disable protrusion for tt fonts
}{}
\makeatletter
\@ifundefined{KOMAClassName}{% if non-KOMA class
  \IfFileExists{parskip.sty}{%
    \usepackage{parskip}
  }{% else
    \setlength{\parindent}{0pt}
    \setlength{\parskip}{6pt plus 2pt minus 1pt}}
}{% if KOMA class
  \KOMAoptions{parskip=half}}
\makeatother
\usepackage{xcolor}
\usepackage[margin=0.7in,footskip=0.4in]{geometry}
\usepackage{color}
\usepackage{fancyvrb}
\newcommand{\VerbBar}{|}
\newcommand{\VERB}{\Verb[commandchars=\\\{\}]}
\DefineVerbatimEnvironment{Highlighting}{Verbatim}{commandchars=\\\{\}}
% Add ',fontsize=\small' for more characters per line
\newenvironment{Shaded}{}{}
\newcommand{\AlertTok}[1]{\textcolor[rgb]{1.00,0.00,0.00}{\textbf{#1}}}
\newcommand{\AnnotationTok}[1]{\textcolor[rgb]{0.38,0.63,0.69}{\textbf{\textit{#1}}}}
\newcommand{\AttributeTok}[1]{\textcolor[rgb]{0.49,0.56,0.16}{#1}}
\newcommand{\BaseNTok}[1]{\textcolor[rgb]{0.25,0.63,0.44}{#1}}
\newcommand{\BuiltInTok}[1]{\textcolor[rgb]{0.00,0.50,0.00}{#1}}
\newcommand{\CharTok}[1]{\textcolor[rgb]{0.25,0.44,0.63}{#1}}
\newcommand{\CommentTok}[1]{\textcolor[rgb]{0.38,0.63,0.69}{\textit{#1}}}
\newcommand{\CommentVarTok}[1]{\textcolor[rgb]{0.38,0.63,0.69}{\textbf{\textit{#1}}}}
\newcommand{\ConstantTok}[1]{\textcolor[rgb]{0.53,0.00,0.00}{#1}}
\newcommand{\ControlFlowTok}[1]{\textcolor[rgb]{0.00,0.44,0.13}{\textbf{#1}}}
\newcommand{\DataTypeTok}[1]{\textcolor[rgb]{0.56,0.13,0.00}{#1}}
\newcommand{\DecValTok}[1]{\textcolor[rgb]{0.25,0.63,0.44}{#1}}
\newcommand{\DocumentationTok}[1]{\textcolor[rgb]{0.73,0.13,0.13}{\textit{#1}}}
\newcommand{\ErrorTok}[1]{\textcolor[rgb]{1.00,0.00,0.00}{\textbf{#1}}}
\newcommand{\ExtensionTok}[1]{#1}
\newcommand{\FloatTok}[1]{\textcolor[rgb]{0.25,0.63,0.44}{#1}}
\newcommand{\FunctionTok}[1]{\textcolor[rgb]{0.02,0.16,0.49}{#1}}
\newcommand{\ImportTok}[1]{\textcolor[rgb]{0.00,0.50,0.00}{\textbf{#1}}}
\newcommand{\InformationTok}[1]{\textcolor[rgb]{0.38,0.63,0.69}{\textbf{\textit{#1}}}}
\newcommand{\KeywordTok}[1]{\textcolor[rgb]{0.00,0.44,0.13}{\textbf{#1}}}
\newcommand{\NormalTok}[1]{#1}
\newcommand{\OperatorTok}[1]{\textcolor[rgb]{0.40,0.40,0.40}{#1}}
\newcommand{\OtherTok}[1]{\textcolor[rgb]{0.00,0.44,0.13}{#1}}
\newcommand{\PreprocessorTok}[1]{\textcolor[rgb]{0.74,0.48,0.00}{#1}}
\newcommand{\RegionMarkerTok}[1]{#1}
\newcommand{\SpecialCharTok}[1]{\textcolor[rgb]{0.25,0.44,0.63}{#1}}
\newcommand{\SpecialStringTok}[1]{\textcolor[rgb]{0.73,0.40,0.53}{#1}}
\newcommand{\StringTok}[1]{\textcolor[rgb]{0.25,0.44,0.63}{#1}}
\newcommand{\VariableTok}[1]{\textcolor[rgb]{0.10,0.09,0.49}{#1}}
\newcommand{\VerbatimStringTok}[1]{\textcolor[rgb]{0.25,0.44,0.63}{#1}}
\newcommand{\WarningTok}[1]{\textcolor[rgb]{0.38,0.63,0.69}{\textbf{\textit{#1}}}}
\usepackage{longtable,booktabs,array}
\usepackage{calc} % for calculating minipage widths
% Correct order of tables after \paragraph or \subparagraph
\usepackage{etoolbox}
\makeatletter
\patchcmd\longtable{\par}{\if@noskipsec\mbox{}\fi\par}{}{}
\makeatother
% Allow footnotes in longtable head/foot
\IfFileExists{footnotehyper.sty}{\usepackage{footnotehyper}}{\usepackage{footnote}}
\makesavenoteenv{longtable}
\usepackage{graphicx}
\makeatletter
\def\maxwidth{\ifdim\Gin@nat@width>\linewidth\linewidth\else\Gin@nat@width\fi}
\def\maxheight{\ifdim\Gin@nat@height>\textheight\textheight\else\Gin@nat@height\fi}
\makeatother
% Scale images if necessary, so that they will not overflow the page
% margins by default, and it is still possible to overwrite the defaults
% using explicit options in \includegraphics[width, height, ...]{}
\setkeys{Gin}{width=\maxwidth,height=\maxheight,keepaspectratio}
% Set default figure placement to htbp
\makeatletter
\def\fps@figure{htbp}
\makeatother
\setlength{\emergencystretch}{3em} % prevent overfull lines
\providecommand{\tightlist}{%
  \setlength{\itemsep}{0pt}\setlength{\parskip}{0pt}}
\setcounter{secnumdepth}{-\maxdimen} % remove section numbering

\usepackage{fancyhdr}
\usepackage{lastpage}
\usepackage{etoolbox}
\newcommand{\tpfoot}{\fancyhf{}\renewcommand{\headrulewidth}{0pt}\renewcommand{\footrulewidth}{0.4pt}%
\fancyfoot[L]{\footnotesize TEST\_PLAN --- SA400 LED Monitor Test Plan · Version 1.1}%
\fancyfoot[R]{\footnotesize Created 2026-10-19 · Updated 2026-10-19 · Page \thepage\ of \pageref*{LastPage}}}
\pagestyle{fancy}\tpfoot
\fancypagestyle{plain}{\tpfoot}
\AtBeginEnvironment{longtable}{\footnotesize}
\AtBeginEnvironment{Shaded}{\footnotesize}
\setlength{\LTpre}{6pt}\setlength{\LTpost}{6pt}
\renewcommand{\arraystretch}{1.25}
\ifLuaTeX
  \usepackage{selnolig}  % disable illegal ligatures
\fi
\IfFileExists{bookmark.sty}{\usepackage{bookmark}}{\usepackage{hyperref}}
\IfFileExists{xurl.sty}{\usepackage{xurl}}{} % add URL line breaks if available
\urlstyle{same}
\hypersetup{
  colorlinks=true,
  linkcolor={blue},
  filecolor={Maroon},
  citecolor={Blue},
  urlcolor={blue},
  pdfcreator={LaTeX via pandoc}}

\author{}
\date{}

\begin{document}

\hypertarget{test-plan-sa400-drive-select-write-gate-led-monitor}{%
\section{Test Plan --- SA400 Drive-Select / Write-Gate LED
Monitor}\label{test-plan-sa400-drive-select-write-gate-led-monitor}}

Document: \texttt{TEST\_PLAN.md} · Project:
\texttt{sa400\_led\_monitor} · Test-plan version
\textbf{1.1} · created 2026-10-19 · updated 2026-10-19 (1.1:
populated-board drawing and legend, §0.7)

Applies to: board rev \textbf{v0.2.2} copper (the 0.2.3
gerber changed Q1 silkscreen only), populated by JLCPCB from
the \textbf{r2} BOM/CPL in \texttt{jlcpcb\_v0.2.2/}.

\begin{center}\rule{0.5\linewidth}{0.5pt}\end{center}

\hypertarget{read-this-first}{%
\subsection{0. Read this first}\label{read-this-first}}

\hypertarget{what-the-three-phases-prove}{%
\subsubsection{0.1 What the three phases
prove}\label{what-the-three-phases-prove}}

\begin{longtable}[]{@{}
  >{\raggedright\arraybackslash}p{(\columnwidth - 6\tabcolsep) * \real{0.2500}}
  >{\raggedright\arraybackslash}p{(\columnwidth - 6\tabcolsep) * \real{0.2500}}
  >{\raggedright\arraybackslash}p{(\columnwidth - 6\tabcolsep) * \real{0.2500}}
  >{\raggedright\arraybackslash}p{(\columnwidth - 6\tabcolsep) * \real{0.2500}}@{}}
\toprule\noalign{}
\begin{minipage}[b]{\linewidth}\raggedright
Phase
\end{minipage} & \begin{minipage}[b]{\linewidth}\raggedright
Setup
\end{minipage} & \begin{minipage}[b]{\linewidth}\raggedright
What it proves
\end{minipage} & \begin{minipage}[b]{\linewidth}\raggedright
Can it damage the TRS-80?
\end{minipage} \\
\midrule\noalign{}
\endhead
\bottomrule\noalign{}
\endlastfoot
\textbf{1 --- Unpowered} & Board only, DMM only & Assembly
is right (orientation, values, no shorts); the bus fingers
connect only where the design says & No --- nothing is
powered \\
\textbf{2 --- Bench power} & Board + current-limited bench
supply + a simple bus-simulator jig & Logic and LEDs work;
the protection works; the \textbf{most} current and voltage
the board can ever put on a bus line, measured & No --- the
EI and drives are not connected \\
\textbf{3 --- On the EI} & Board on the real cable with the
Expansion Interface, then with drives & Works with real
signals; does not disturb disk I/O; safe in every
power-on/power-off order & Only if a Phase 1 or 2 check was
skipped or failed \\
\end{longtable}

\hypertarget{rules}{%
\subsubsection{0.2 Rules}\label{rules}}

\begin{enumerate}
\def\labelenumi{\arabic{enumi}.}
\tightlist
\item
  \textbf{Do the phases in order, on each board.} A board
  does not go to the next phase until every check in the
  current phase passes.
\item
  \textbf{Stop on the first out-of-spec reading.} Power off,
  write it down, and fix or diagnose before going on. Do not
  ``see if it works anyway'' with the EI attached.
\item
  \textbf{Never plug or unplug the card edge (J1) with
  anything powered.} Power (J2) may be unplugged while
  running only in step 3.5, which is designed for it.
\item
  \textbf{Board 1 goes all the way through first.} Boards
  2--5 do Phase 1 and Phase 2 sections 2.1--2.6. Only board
  1 does the protection tests in 2.7 and 2.8 (they stress
  parts on purpose).
\item
  Write a serial number (1--5) on each board's silkscreen
  with a marker, and record results in the sheets in section
  5.
\end{enumerate}

\hypertarget{as-built-values-r2-bom-these-differ-from-older-docs}{%
\subsubsection{0.3 As-built values (r2 BOM): these differ
from older
docs}\label{as-built-values-r2-bom-these-differ-from-older-docs}}

\texttt{docs/BOM.md} and \texttt{THEORY\_OF\_OPERATION.md}
still show the design-time values. The boards were built
from the r2 BOM, so \textbf{this plan uses the r2 values}:

\begin{longtable}[]{@{}
  >{\raggedright\arraybackslash}p{(\columnwidth - 6\tabcolsep) * \real{0.2500}}
  >{\raggedright\arraybackslash}p{(\columnwidth - 6\tabcolsep) * \real{0.2500}}
  >{\raggedright\arraybackslash}p{(\columnwidth - 6\tabcolsep) * \real{0.2500}}
  >{\raggedright\arraybackslash}p{(\columnwidth - 6\tabcolsep) * \real{0.2500}}@{}}
\toprule\noalign{}
\begin{minipage}[b]{\linewidth}\raggedright
Ref
\end{minipage} & \begin{minipage}[b]{\linewidth}\raggedright
Design docs say
\end{minipage} & \begin{minipage}[b]{\linewidth}\raggedright
\textbf{As built (r2)}
\end{minipage} & \begin{minipage}[b]{\linewidth}\raggedright
Effect on readings
\end{minipage} \\
\midrule\noalign{}
\endhead
\bottomrule\noalign{}
\endlastfoot
R1--R8 (drive LEDs), R21 (5V LED) & 1.5 kΩ & \textbf{470 Ω}
& LEDs run ≈ 5.5--6.5 mA (still inside the 74LS00's 8 mA
sink rating) \\
R19 (12V LED) & 4.7 kΩ & \textbf{2.2 kΩ} & D12 runs ≈ 4.5
mA \\
Q1 & NDP6020P & \textbf{IRF4905PBF} (same G-D-S pinout) &
Specified for V\_GS = −10 V; at our −5 V it should still be
fully on at 35 mA, and \textbf{step 2.3 measures that} \\
R9--R13, R14--R18 & 10 kΩ, 470 Ω & same & --- \\
\end{longtable}

\hypertarget{reference-map-used-throughout}{%
\subsubsection{0.4 Reference map (used
throughout)}\label{reference-map-used-throughout}}

\textbf{J1 card edge.} Even pins are on the component side,
odd pins on the solder side directly underneath.

\begin{longtable}[]{@{}
  >{\raggedright\arraybackslash}p{(\columnwidth - 16\tabcolsep) * \real{0.1111}}
  >{\raggedright\arraybackslash}p{(\columnwidth - 16\tabcolsep) * \real{0.1111}}
  >{\raggedright\arraybackslash}p{(\columnwidth - 16\tabcolsep) * \real{0.1111}}
  >{\raggedright\arraybackslash}p{(\columnwidth - 16\tabcolsep) * \real{0.1111}}
  >{\raggedright\arraybackslash}p{(\columnwidth - 16\tabcolsep) * \real{0.1111}}
  >{\raggedright\arraybackslash}p{(\columnwidth - 16\tabcolsep) * \real{0.1111}}
  >{\raggedright\arraybackslash}p{(\columnwidth - 16\tabcolsep) * \real{0.1111}}
  >{\raggedright\arraybackslash}p{(\columnwidth - 16\tabcolsep) * \real{0.1111}}
  >{\raggedright\arraybackslash}p{(\columnwidth - 16\tabcolsep) * \real{0.1111}}@{}}
\toprule\noalign{}
\begin{minipage}[b]{\linewidth}\raggedright
Signal
\end{minipage} & \begin{minipage}[b]{\linewidth}\raggedright
J1 pin
\end{minipage} & \begin{minipage}[b]{\linewidth}\raggedright
Its GND partner (under it)
\end{minipage} & \begin{minipage}[b]{\linewidth}\raggedright
Series R (finger → receiver)
\end{minipage} & \begin{minipage}[b]{\linewidth}\raggedright
Pull-up
\end{minipage} & \begin{minipage}[b]{\linewidth}\raggedright
U1 input
\end{minipage} & \begin{minipage}[b]{\linewidth}\raggedright
U1 output
\end{minipage} & \begin{minipage}[b]{\linewidth}\raggedright
Red LED / driver
\end{minipage} & \begin{minipage}[b]{\linewidth}\raggedright
Green LED / driver
\end{minipage} \\
\midrule\noalign{}
\endhead
\bottomrule\noalign{}
\endlastfoot
DS1 (TRS-80 drive :0) & 10 & 9 & R17 & R9 & U1-1 & U1-2 & D1
/ U3-2 & D5 / U2-3 \\
DS2 (drive :1) & 12 & 11 & R16 & R10 & U1-3 & U1-4 & D2 /
U3-4 & D6 / U2-6 \\
DS3 (drive :2) & 14 & 13 & R15 & R11 & U1-5 & U1-6 & D3 /
U3-6 & D7 / U2-8 \\
DS4 (drive :3) & 6 & 5 & R18 & R12 & U1-9 & U1-8 & D4 / U3-8
& D8 / U2-11 \\
WRITE GATE & 24 & 23 & R14 & R13 & U1-11 & U1-10 & --- &
(all greens through U2) \\
\end{longtable}

\begin{itemize}
\tightlist
\item
  \textbf{GND fingers (15):} 3, 5, 7, 9, 11, 13, 15, 17, 19,
  21, 23, 25, 27, 29, 31.
\item
  \textbf{No-connect fingers (14):} 1 and 33 (they sit over
  the clip holes), and 2, 4, 8, 16, 18, 20, 22, 26, 28, 30,
  32, 34.
\end{itemize}

\textbf{J2 power (Mate-N-Lok, solder side, pin 1 = square
pad).} 1 = +12 V, 2 = +12 RTN, 3 = +5 RTN, 4 = +5 V. This is
the same connector and pinout as a PC ``Molex'' 4-pin
peripheral plug (yellow +12, black, black, red +5).

\textbf{Power path.} J2-4 → F1 (PTC) → Q1 pin 2 (drain) → Q1
pin 3 (source) = \textbf{+5V rail} (also U1/U2/U3 pin 14,
D11 cathode, C4+). Q1 pin 1 (gate) = GND. J2-1 → R19 → D12 →
GND only.

\textbf{Handy test points.} The vertical resistors stand up
with one lead looped over the top. That loop is a good probe
point, but which net it is on depends on which way JLCPCB
stood the part up. \textbf{Before you use any resistor loop
as a test point, check it with the DMM against the finger or
the IC pin.}

\hypertarget{equipment}{%
\subsubsection{0.5 Equipment}\label{equipment}}

\begin{longtable}[]{@{}
  >{\raggedright\arraybackslash}p{(\columnwidth - 4\tabcolsep) * \real{0.3333}}
  >{\raggedright\arraybackslash}p{(\columnwidth - 4\tabcolsep) * \real{0.3333}}
  >{\raggedright\arraybackslash}p{(\columnwidth - 4\tabcolsep) * \real{0.3333}}@{}}
\toprule\noalign{}
\begin{minipage}[b]{\linewidth}\raggedright
Item
\end{minipage} & \begin{minipage}[b]{\linewidth}\raggedright
Needed for
\end{minipage} & \begin{minipage}[b]{\linewidth}\raggedright
Notes
\end{minipage} \\
\midrule\noalign{}
\endhead
\bottomrule\noalign{}
\endlastfoot
DMM with ohms, diode, DC volts, DC mA & all & A second DMM
to watch current while you probe saves time \\
Bench supply, adjustable 0--15 V, \textbf{adjustable current
limit} & 2, 3 & Two outputs is ideal (one for the board, one
for the jig pull-ups). The outputs must be floating (not
tied to earth on +) \\
Power pigtail for J2 & 2, 3 & TE/AMP 1-480424-0 housing +
60619-1 pins ---
https://www.te.com/en/product-1-480424-0.html · pins
https://www.mouser.com/c/?q=60619-1 --- or any PC
Molex-to-bare-wire cable (same connector). Label the
wires. \\
\textbf{Bus-simulator jig} (section 0.6) & 2 & Recommended.
The fallback is clip leads on the fingers, which risks
scratching the gold or bridging two fingers \\
IC extractor, ESD strap & 1, 2.7, 2.8 & ICs are socketed, so
you can pull them for the protection tests \\
Scope (optional) & 3.3 & Best way to show the board doesn't
change bus levels or edges \\
IR thermometer or fingertip & 2, 3.6 & Heat check \\
Scratch diskettes (2+) & 3 & \textbf{Never test writes on a
disk you care about} \\
\end{longtable}

\hypertarget{bus-simulator-jig-fake-ei}{%
\subsubsection{0.6 Bus-simulator jig (``fake
EI'')}\label{bus-simulator-jig-fake-ei}}

This stands in for the EI's open-collector drivers and the
terminating drive. Each line is held high by 150 Ω (like the
terminator), and a switch pulls it to ground (like a 7438
asserting the line).

\begin{verbatim}
 3M 3463-0001 card-edge IDC  --  ~150 mm 34-way ribbon  --  cut end; separate:
   (red stripe = pin 1)            signals : 6, 10, 12, 14, 24
                                   GND     : 5, 9, 11, 13, 23 (twist together)

 Each signal conductor (x5):

   +5V_JIG --[JP on/off]--[150 R]--+-- conductor 6 / 10 / 12 / 14 / 24
                                   |
                                 [SW n]   closed = ASSERTED
                                   |
                                  GND --- GND conductors + bench supply (-)
\end{verbatim}

\begin{itemize}
\tightlist
\item
  Use a 3463-0001 from the cable build:
  https://www.mouser.com/c/?q=3463-0001
\item
  Ribbon conductor N = J1 pin N when the stripe is on pin 1
  of the IDC connector. Before you trust the jig, check that
  with the DMM.
\item
  \textbf{JP} turns the five 150 Ω pull-ups on or off.
  \emph{Off} = the board's own 10 k pull-ups only. \emph{On}
  = simulates a terminated bus. +5V\_JIG comes from the
  second bench output (or the same output, but then the
  board-current reading also includes the jig current).
\item
  Five switches, labelled DS1 DS2 DS3 DS4 WG.
\end{itemize}

\clearpage
\hypertarget{populated-board-wireframe-and-legend}{%
\subsubsection{0.7 Populated board (wireframe) and
legend}\label{populated-board-wireframe-and-legend}}

\begin{figure}
\centering
\includegraphics[width=1\textwidth,height=\textheight]{img/board_populated.pdf}
\caption{Populated board, component side, wireframe ---
generated by \texttt{gen/render\_populated.py} from the
footprint positions in \texttt{gen/design.py}, as-built
values}
\end{figure}

\clearpage
\begin{longtable}[]{@{}
  >{\raggedright\arraybackslash}p{(\columnwidth - 8\tabcolsep) * \real{0.2000}}
  >{\raggedright\arraybackslash}p{(\columnwidth - 8\tabcolsep) * \real{0.2000}}
  >{\raggedright\arraybackslash}p{(\columnwidth - 8\tabcolsep) * \real{0.2000}}
  >{\raggedright\arraybackslash}p{(\columnwidth - 8\tabcolsep) * \real{0.2000}}
  >{\raggedright\arraybackslash}p{(\columnwidth - 8\tabcolsep) * \real{0.2000}}@{}}
\toprule\noalign{}
\begin{minipage}[b]{\linewidth}\raggedright
Ref
\end{minipage} & \begin{minipage}[b]{\linewidth}\raggedright
Part (as built, r2 BOM)
\end{minipage} & \begin{minipage}[b]{\linewidth}\raggedright
Drawn as
\end{minipage} & \begin{minipage}[b]{\linewidth}\raggedright
Job on the board
\end{minipage} & \begin{minipage}[b]{\linewidth}\raggedright
Check when inspecting
\end{minipage} \\
\midrule\noalign{}
\endhead
\bottomrule\noalign{}
\endlastfoot
U1 & 74LS14 hex Schmitt inverter, DIP-14, in socket & Solid
rounded box with end notch, inside a dashed socket outline &
Bus receivers for DS1--DS4 and WRITE GATE & Notch at left;
pin 1 = square pad, bottom-left \\
U2 & 74LS00 quad NAND, DIP-14, in socket & Same & Green
LEDs: DSn AND WG & Same \\
U3 & 74LS06 hex open-collector inverter, DIP-14, in socket &
Same & Red LEDs: DSn & Same \\
Q1 & IRF4905 P-channel MOSFET, TO-220, standing & Rectangle
(body) + thin rectangle (metal tab), pins G D S &
Reverse-polarity switch on +5 V & Gate in the square pad;
tab toward the front edge \\
F1 & PTC resettable fuse 0.5 A & Rounded bar & +5 V input
fuse & --- \\
D11 & 1N5339B 5.6 V 5 W Zener, DO-15, standing & Double
circle (body + band) with the lead down to the anode & +5 V
rail clamp & Band (cathode) on top, in the square pad \\
C4 & 10 µF 16 V electrolytic & Circle with ``+'' and a line
on the stripe side & Bulk +5 V & + lead in the square pad \\
C1, C2, C3, C6 & 100 nF ceramic disc & Small rounded slot &
Decoupling U1/U2/U3; rail entry (C6) & --- \\
R14--R18 & 470 Ω ¼ W, standing & Circle (body) + line (lead)
& Bus series resistors, finger → receiver (fault-current
limit) & Value printed under the lead \\
R9--R13 & 10 kΩ ¼ W, standing & Same, label rotated & Bus
pull-ups on the receiver side & --- \\
R1--R8 & 470 Ω ¼ W, standing & Same & Drive-LED series
resistors (R1--R4 red, R5--R8 green) & --- \\
R19 & 2.2 kΩ ¼ W & Same & +12 V indicator series & --- \\
R21 & 470 Ω ¼ W & Same & +5 V indicator series & --- \\
D1--D4 & Red 5 mm LED (R) & Circle with a flat; ``K'' at the
flat & Drive 1--4 SELECTED & Flat / short lead = cathode =
square pad \\
D5--D8 & Green 5 mm LED (G) & Same & Drive 1--4 WRITE GATE &
Same \\
D12, D13 & Yellow 5 mm LED (Y) & Same & ``12V'' and ``5V''
present & Same \\
J1 & 34-contact card edge (gold fingers) & Rectangles on the
tab; bold numbers = monitored lines 6/10/12/14/24 & Shugart
bus & Even pins on top, odd under them; key slot between 4
and 6 \\
J2 & TE/AMP Mate-N-Lok 350211-1 & \textbf{Dashed grey}: it
sits on the solder side & Power in: 1 +12 V, 2/3 RTN, 4 +5 V
& Pin 1 = square pad \\
H1--H4 & M3 mounting holes & Circle in a dashed ring &
Standoffs & --- \\
\end{longtable}

Conventions: a \textbf{square pad} is pin 1, a cathode, or a
capacitor's + lead. \textbf{Dashed} means a socket or a part
on the solder side. Drive \emph{n} on the board is TRS-80
drive \emph{:n−1}.

\begin{center}\rule{0.5\linewidth}{0.5pt}\end{center}

\clearpage
\hypertarget{phase-1-unpowered-no-supply-not-on-the-ei}{%
\subsection{Phase 1 --- Unpowered (no supply, not on the
EI)}\label{phase-1-unpowered-no-supply-not-on-the-ei}}

Board fully assembled, ICs in their sockets, nothing plugged
into J1 or J2. Hold the DMM leads still while a reading
settles, because the capacitors charge through the meter.

\hypertarget{visual-and-mechanical}{%
\subsubsection{1.1 Visual and
mechanical}\label{visual-and-mechanical}}

\begin{longtable}[]{@{}
  >{\raggedright\arraybackslash}p{(\columnwidth - 4\tabcolsep) * \real{0.3333}}
  >{\raggedright\arraybackslash}p{(\columnwidth - 4\tabcolsep) * \real{0.3333}}
  >{\raggedright\arraybackslash}p{(\columnwidth - 4\tabcolsep) * \real{0.3333}}@{}}
\toprule\noalign{}
\begin{minipage}[b]{\linewidth}\raggedright
\#
\end{minipage} & \begin{minipage}[b]{\linewidth}\raggedright
Check
\end{minipage} & \begin{minipage}[b]{\linewidth}\raggedright
Pass
\end{minipage} \\
\midrule\noalign{}
\endhead
\bottomrule\noalign{}
\endlastfoot
1.1.1 & U1 = 74LS14, U2 = 74LS00, U3 = 74LS06, each in the
correct position, notch matches the silk, every pin inside
the socket (none folded under) & ✓ \\
1.1.2 & Q1 is the IRF4905: pin 1 (gate) in the square pad,
metal tab on the side shown by the silkscreen & ✓ \\
1.1.3 & D11: band (cathode) at the square pad. C4: + lead at
the square pad & ✓ \\
1.1.4 & All LEDs: flat side / short lead (cathode) at the
square pad. Colours: D1--D4 red, D5--D8 green, D12--D13
yellow & ✓ \\
1.1.5 & Resistor colour bands: R1--R8, R21, R14--R18 = 470 Ω
(yellow-violet-brown); R9--R13 = 10 k (brown-black-orange);
R19 = 2.2 k (red-red-black-brown for 1 \%) & ✓ \\
1.1.6 & J2 is on the \textbf{solder side}, housing opening
toward the rear edge, seated flat & ✓ \\
1.1.7 & Solder side: no bridges, especially across the J2
pins and the DIP rows. No flux balls & ✓ \\
1.1.8 & Fingers: gold, bevelled, \textbf{no solder or flux
on any finger}, no lifted finger & ✓ \\
1.1.9 & Board thickness at the fingers (calipers):
1.37--1.78 mm (3M spec), expect 1.6 & record \\
1.1.10 & Fit test: push a \textbf{bare} 3463 connector (not
on a live cable) onto the tab. It seats fully, and the key
(if fitted) goes into the slot between pins 4 and 6. Remove
it & ✓ \\
1.1.11 & Fit test: the power plug latches into J2 and only
goes in one way & ✓ \\
1.1.12 & \textbf{Record the three ASSUMED dimensions} while
you have calipers out: tab edge to corner (5.08 assumed),
key-plug width (0.91), J2 pin drill (2.1) --- see
\texttt{docs/MECHANICAL.md} & record \\
\end{longtable}

\hypertarget{power-connector-no-shorts-between-rails}{%
\subsubsection{1.2 Power connector --- no shorts between
rails}\label{power-connector-no-shorts-between-rails}}

DMM in ohms unless stated. ``OL'' = open / over-range.

\begin{longtable}[]{@{}
  >{\raggedright\arraybackslash}p{(\columnwidth - 8\tabcolsep) * \real{0.2000}}
  >{\raggedright\arraybackslash}p{(\columnwidth - 8\tabcolsep) * \real{0.2000}}
  >{\raggedright\arraybackslash}p{(\columnwidth - 8\tabcolsep) * \real{0.2000}}
  >{\raggedright\arraybackslash}p{(\columnwidth - 8\tabcolsep) * \real{0.2000}}
  >{\raggedright\arraybackslash}p{(\columnwidth - 8\tabcolsep) * \real{0.2000}}@{}}
\toprule\noalign{}
\begin{minipage}[b]{\linewidth}\raggedright
\#
\end{minipage} & \begin{minipage}[b]{\linewidth}\raggedright
From (red)
\end{minipage} & \begin{minipage}[b]{\linewidth}\raggedright
To (black)
\end{minipage} & \begin{minipage}[b]{\linewidth}\raggedright
Expect
\end{minipage} & \begin{minipage}[b]{\linewidth}\raggedright
Proves
\end{minipage} \\
\midrule\noalign{}
\endhead
\bottomrule\noalign{}
\endlastfoot
1.2.1 & J2-1 (+12) & J2-4 (+5) & \textbf{OL}, then try the
leads the other way: \textbf{OL} & 12 V can never reach the
5 V rail through the board \\
1.2.2 & J2-2 & J2-3 & \textless{} 1 Ω & Both returns are
tied together \\
1.2.3 & J2-3 & J1-9 (any GND finger) & \textless{} 1 Ω &
Supply ground = bus ground \\
1.2.4 & J2-4 & J2-3 (GND) & \textgreater{} 200 Ω and
settling (not a dead short) & No short across the +5 V
input \\
1.2.5 & J2-1 & J2-3 & \textgreater{} 2 kΩ or OL & No short
across the +12 V input \\
1.2.6 & J2-4 & F1 far lead / Q1 pin 2 & \textless{} 2 Ω
(cold PTC) & F1 is in circuit and intact \\
1.2.7 & Q1 pin 1 (gate) & GND & \textless{} 1 Ω & Gate tied
to ground as designed \\
\end{longtable}

\hypertarget{reverse-polarity-switch-q1-orientation-diode-mode}{%
\subsubsection{1.3 Reverse-polarity switch (Q1) orientation
--- diode
mode}\label{reverse-polarity-switch-q1-orientation-diode-mode}}

\begin{longtable}[]{@{}
  >{\raggedright\arraybackslash}p{(\columnwidth - 8\tabcolsep) * \real{0.2000}}
  >{\raggedright\arraybackslash}p{(\columnwidth - 8\tabcolsep) * \real{0.2000}}
  >{\raggedright\arraybackslash}p{(\columnwidth - 8\tabcolsep) * \real{0.2000}}
  >{\raggedright\arraybackslash}p{(\columnwidth - 8\tabcolsep) * \real{0.2000}}
  >{\raggedright\arraybackslash}p{(\columnwidth - 8\tabcolsep) * \real{0.2000}}@{}}
\toprule\noalign{}
\begin{minipage}[b]{\linewidth}\raggedright
\#
\end{minipage} & \begin{minipage}[b]{\linewidth}\raggedright
Red lead
\end{minipage} & \begin{minipage}[b]{\linewidth}\raggedright
Black lead
\end{minipage} & \begin{minipage}[b]{\linewidth}\raggedright
Expect
\end{minipage} & \begin{minipage}[b]{\linewidth}\raggedright
Proves
\end{minipage} \\
\midrule\noalign{}
\endhead
\bottomrule\noalign{}
\endlastfoot
1.3.1 & Q1 pin 2 (drain) & Q1 pin 3 (source / rail) &
0.4--0.7 V & Body diode points from input to rail, so Q1 is
fitted the right way round \\
1.3.2 & Q1 pin 3 & Q1 pin 2 & OL (may start low and climb
while C4 charges) & Reversed supply will be blocked \\
1.3.3 & J2-4 & U1-14 (+5V rail) & 0.4--0.7 V & Whole input
path F1 → Q1 → rail is in one piece \\
1.3.4 & GND & U1-14 & 0.4--0.8 V & Normal: D11 and the IC
substrate diodes. \textbf{≈ 0 V = short on the rail} \\
\end{longtable}

If 1.3.1 reads OL and 1.3.2 reads 0.5 V, Q1 is in backwards.
\textbf{Stop.}

\hypertarget{bus-fingers-the-board-touches-only-what-it-should}{%
\subsubsection{1.4 Bus fingers --- the board touches only
what it
should}\label{bus-fingers-the-board-touches-only-what-it-should}}

This is the core ``cannot hurt the EI'' check with no power
applied.

\begin{longtable}[]{@{}
  >{\raggedright\arraybackslash}p{(\columnwidth - 4\tabcolsep) * \real{0.3333}}
  >{\raggedright\arraybackslash}p{(\columnwidth - 4\tabcolsep) * \real{0.3333}}
  >{\raggedright\arraybackslash}p{(\columnwidth - 4\tabcolsep) * \real{0.3333}}@{}}
\toprule\noalign{}
\begin{minipage}[b]{\linewidth}\raggedright
\#
\end{minipage} & \begin{minipage}[b]{\linewidth}\raggedright
Test
\end{minipage} & \begin{minipage}[b]{\linewidth}\raggedright
Expect
\end{minipage} \\
\midrule\noalign{}
\endhead
\bottomrule\noalign{}
\endlastfoot
1.4.1 & Each of the \textbf{14 no-connect fingers} (1, 33,
2, 4, 8, 16, 18, 20, 22, 26, 28, 30, 32, 34) to GND (J2-3) &
\textbf{OL} \\
1.4.2 & The same 14 fingers to the +5V rail (U1-14) &
\textbf{OL} \\
1.4.3 & Each of the \textbf{15 GND fingers} to J2-3 &
\textless{} 1 Ω \\
1.4.4 & Each monitored finger to its U1 input (table 0.4,
e.g.~J1-10 → U1-1) & \textbf{447--494 Ω} (470 Ω ±5 \%) \\
1.4.5 & Each monitored finger to the +5V rail, red lead on
the finger & \textbf{9.9--11.0 kΩ} (470 Ω + 10 kΩ) \\
1.4.6 & Each monitored finger to GND & \textgreater{} 10 kΩ,
\textbf{never \textless{} 1 kΩ} \\
1.4.7 & Adjacent signal fingers to each other: 10↔12, 12↔14,
4↔6, 6↔8, 22↔24, 24↔26 & ≈ 21 kΩ between two monitored
fingers (each one's 10.47 k to the rail), \textbf{OL} when
one of the pair is a no-connect \\
1.4.8 & Each monitored finger and J2-1 (+12) & \textbf{OL}
--- 12 V has no path to the bus \\
\end{longtable}

\hypertarget{leds-in-circuit-diode-mode}{%
\subsubsection{1.5 LEDs in circuit --- diode
mode}\label{leds-in-circuit-diode-mode}}

Red lead on the anode (round pad), black lead on the cathode
(square pad). Most DMMs light a red or yellow LED faintly.
Some can't light green.

\begin{longtable}[]{@{}
  >{\raggedright\arraybackslash}p{(\columnwidth - 4\tabcolsep) * \real{0.3333}}
  >{\raggedright\arraybackslash}p{(\columnwidth - 4\tabcolsep) * \real{0.3333}}
  >{\raggedright\arraybackslash}p{(\columnwidth - 4\tabcolsep) * \real{0.3333}}@{}}
\toprule\noalign{}
\begin{minipage}[b]{\linewidth}\raggedright
\#
\end{minipage} & \begin{minipage}[b]{\linewidth}\raggedright
LED
\end{minipage} & \begin{minipage}[b]{\linewidth}\raggedright
Expect
\end{minipage} \\
\midrule\noalign{}
\endhead
\bottomrule\noalign{}
\endlastfoot
1.5.1 & D1--D4, D12, D13 & 1.5--2.1 V, faint glow \\
1.5.2 & D5--D8 & 1.7--2.4 V or OL (meter can't reach a
green's Vf); faint glow if it can \\
1.5.3 & Each LED reversed & OL or a much higher reading
(other in-circuit paths) \\
\end{longtable}

Phase 1 done → record, go to Phase 2.

\begin{center}\rule{0.5\linewidth}{0.5pt}\end{center}

\hypertarget{phase-2-bench-power-not-on-the-ei}{%
\subsection{Phase 2 --- Bench power, not on the
EI}\label{phase-2-bench-power-not-on-the-ei}}

Board on the bench, J2 fed from the bench supply through the
pigtail. J1 is either empty or plugged into the \textbf{jig
only}. The EI, drives and cable are nowhere near it.

Expected currents (as built):

\begin{longtable}[]{@{}
  >{\raggedright\arraybackslash}p{(\columnwidth - 2\tabcolsep) * \real{0.5000}}
  >{\raggedright\arraybackslash}p{(\columnwidth - 2\tabcolsep) * \real{0.5000}}@{}}
\toprule\noalign{}
\begin{minipage}[b]{\linewidth}\raggedright
Condition
\end{minipage} & \begin{minipage}[b]{\linewidth}\raggedright
Expected +5 V current
\end{minipage} \\
\midrule\noalign{}
\endhead
\bottomrule\noalign{}
\endlastfoot
Idle (no line asserted) & \textbf{≈ 25--45 mA} (three LS ICs
+ D13 at ≈ 6 mA) --- accept 15--60 mA, but boards should
agree within ±20 \% \\
Each asserted line (jig pull-ups off) & + ≈ 0.48 mA (its
10.47 k pull-up path) + any LED it lights \\
One drive selected & + ≈ 5.5--6 mA (red) \\
Selected and writing & + another ≈ 5.5 mA (green) \\
+12 V input & ≈ 4.5 mA (D12 only) \\
\end{longtable}

\hypertarget{first-power-up-5-v-only-current-limited}{%
\subsubsection{2.1 First power-up --- +5 V only,
current-limited}\label{first-power-up-5-v-only-current-limited}}

\begin{enumerate}
\def\labelenumi{\arabic{enumi}.}
\tightlist
\item
  Supply output \textbf{off}. Set it to \textbf{0 V, 100 mA
  limit}. Connect J2-4 (+) and J2-3 (−). Leave J2-1 open.
  Nothing on J1.
\item
  Turn the output on and \textbf{raise slowly} to 5.00 V,
  watching the current.

  \begin{itemize}
  \tightlist
  \item
    Over 60 mA at any point, or the supply drops into
    current limit → \textbf{stop}: short, or a reversed IC.
  \end{itemize}
\item
  At 5.00 V record:
\end{enumerate}

\begin{longtable}[]{@{}
  >{\raggedright\arraybackslash}p{(\columnwidth - 4\tabcolsep) * \real{0.3333}}
  >{\raggedright\arraybackslash}p{(\columnwidth - 4\tabcolsep) * \real{0.3333}}
  >{\raggedright\arraybackslash}p{(\columnwidth - 4\tabcolsep) * \real{0.3333}}@{}}
\toprule\noalign{}
\begin{minipage}[b]{\linewidth}\raggedright
\#
\end{minipage} & \begin{minipage}[b]{\linewidth}\raggedright
Measure
\end{minipage} & \begin{minipage}[b]{\linewidth}\raggedright
Expect
\end{minipage} \\
\midrule\noalign{}
\endhead
\bottomrule\noalign{}
\endlastfoot
2.1.1 & Supply current & 15--60 mA (record exactly --- this
is the board's fingerprint) \\
2.1.2 & D13 ``5V'' & lit \\
2.1.3 & D1--D8, D12 & \textbf{all off} (no line asserted;
D12 has no +12 V) \\
2.1.4 & Touch test after 1 min: U1--U3, Q1, D11, F1 & none
warm \\
\end{longtable}

\hypertarget{rail-voltages}{%
\subsubsection{2.2 Rail voltages}\label{rail-voltages}}

\begin{longtable}[]{@{}
  >{\raggedright\arraybackslash}p{(\columnwidth - 4\tabcolsep) * \real{0.3333}}
  >{\raggedright\arraybackslash}p{(\columnwidth - 4\tabcolsep) * \real{0.3333}}
  >{\raggedright\arraybackslash}p{(\columnwidth - 4\tabcolsep) * \real{0.3333}}@{}}
\toprule\noalign{}
\begin{minipage}[b]{\linewidth}\raggedright
\#
\end{minipage} & \begin{minipage}[b]{\linewidth}\raggedright
Measure (black on GND)
\end{minipage} & \begin{minipage}[b]{\linewidth}\raggedright
Expect
\end{minipage} \\
\midrule\noalign{}
\endhead
\bottomrule\noalign{}
\endlastfoot
2.2.1 & J2-4 & 5.00 V (set) \\
2.2.2 & +5V rail at U1-14, U2-14, U3-14 & within 0.05 V of
J2-4 \\
2.2.3 & U1-7, U2-7, U3-7 & 0.00 V \\
\end{longtable}

\hypertarget{q1-voltage-drop}{%
\subsubsection{2.3 Q1 voltage drop}\label{q1-voltage-drop}}

\begin{longtable}[]{@{}
  >{\raggedright\arraybackslash}p{(\columnwidth - 6\tabcolsep) * \real{0.2500}}
  >{\raggedright\arraybackslash}p{(\columnwidth - 6\tabcolsep) * \real{0.2500}}
  >{\raggedright\arraybackslash}p{(\columnwidth - 6\tabcolsep) * \real{0.2500}}
  >{\raggedright\arraybackslash}p{(\columnwidth - 6\tabcolsep) * \real{0.2500}}@{}}
\toprule\noalign{}
\begin{minipage}[b]{\linewidth}\raggedright
\#
\end{minipage} & \begin{minipage}[b]{\linewidth}\raggedright
Measure
\end{minipage} & \begin{minipage}[b]{\linewidth}\raggedright
Expect
\end{minipage} & \begin{minipage}[b]{\linewidth}\raggedright
If not
\end{minipage} \\
\midrule\noalign{}
\endhead
\bottomrule\noalign{}
\endlastfoot
2.3.1 & DC volts, red J2-4, black U1-14 &
\textbf{\textless{} 0.05 V} & 0.4--0.7 V means Q1's channel
is not turning on and the current is going through its body
diode. The rail would then sit around 4.4 V, below the 74LS
minimum of 4.75 V. Check Q1 orientation and part; the
IRF4905 should be fully on at −5 V V\_GS for 35 mA \\
\end{longtable}

\hypertarget{supply-tolerance-and-clamp-leakage}{%
\subsubsection{2.4 Supply tolerance and clamp
leakage}\label{supply-tolerance-and-clamp-leakage}}

\begin{longtable}[]{@{}
  >{\raggedright\arraybackslash}p{(\columnwidth - 6\tabcolsep) * \real{0.2500}}
  >{\raggedright\arraybackslash}p{(\columnwidth - 6\tabcolsep) * \real{0.2500}}
  >{\raggedright\arraybackslash}p{(\columnwidth - 6\tabcolsep) * \real{0.2500}}
  >{\raggedright\arraybackslash}p{(\columnwidth - 6\tabcolsep) * \real{0.2500}}@{}}
\toprule\noalign{}
\begin{minipage}[b]{\linewidth}\raggedright
\#
\end{minipage} & \begin{minipage}[b]{\linewidth}\raggedright
Set J2-4 to
\end{minipage} & \begin{minipage}[b]{\linewidth}\raggedright
Measure
\end{minipage} & \begin{minipage}[b]{\linewidth}\raggedright
Expect
\end{minipage} \\
\midrule\noalign{}
\endhead
\bottomrule\noalign{}
\endlastfoot
2.4.1 & 4.75 V & current; then repeat 2.5 for one drive &
LEDs still work \\
2.4.2 & 5.25 V & current & No more than \textbf{5 mA} above
the 5.00 V reading. A bigger rise means D11 (5.32 V minimum
knee) is starting to conduct; record it \\
2.4.3 & back to 5.00 V & --- & --- \\
\end{longtable}

\hypertarget{v-indicator}{%
\subsubsection{2.5 +12 V indicator}\label{v-indicator}}

\begin{enumerate}
\def\labelenumi{\arabic{enumi}.}
\tightlist
\item
  Second output (or the same supply in a second step): 12.0
  V, \textbf{50 mA limit}, J2-1 (+) / J2-2 (−).
\end{enumerate}

\begin{longtable}[]{@{}
  >{\raggedright\arraybackslash}p{(\columnwidth - 4\tabcolsep) * \real{0.3333}}
  >{\raggedright\arraybackslash}p{(\columnwidth - 4\tabcolsep) * \real{0.3333}}
  >{\raggedright\arraybackslash}p{(\columnwidth - 4\tabcolsep) * \real{0.3333}}@{}}
\toprule\noalign{}
\begin{minipage}[b]{\linewidth}\raggedright
\#
\end{minipage} & \begin{minipage}[b]{\linewidth}\raggedright
Measure
\end{minipage} & \begin{minipage}[b]{\linewidth}\raggedright
Expect
\end{minipage} \\
\midrule\noalign{}
\endhead
\bottomrule\noalign{}
\endlastfoot
2.5.1 & +12 V current & 3.5--5.5 mA \\
2.5.2 & D12 ``12V'' & lit \\
2.5.3 & +5V rail (U1-14) & unchanged from 2.2 --- 12 V does
not leak onto the logic \\
2.5.4 & Every monitored finger (6, 10, 12, 14, 24), open
circuit & ≈ +5 V, \textbf{never above the rail} \\
\end{longtable}

\hypertarget{bus-isolation-and-logic-with-the-jig}{%
\subsubsection{2.6 Bus isolation and logic --- with the
jig}\label{bus-isolation-and-logic-with-the-jig}}

The monitor stays at 5.00 V / 100 mA limit. Plug the
\textbf{jig} (pull-ups \textbf{off}, all switches open) onto
J1 with the board's supply \textbf{off}, then power up.

\textbf{2.6.A --- What the board puts on the bus (the ``will
it short the EI or a drive'' numbers)}

\begin{longtable}[]{@{}
  >{\raggedright\arraybackslash}p{(\columnwidth - 6\tabcolsep) * \real{0.2500}}
  >{\raggedright\arraybackslash}p{(\columnwidth - 6\tabcolsep) * \real{0.2500}}
  >{\raggedright\arraybackslash}p{(\columnwidth - 6\tabcolsep) * \real{0.2500}}
  >{\raggedright\arraybackslash}p{(\columnwidth - 6\tabcolsep) * \real{0.2500}}@{}}
\toprule\noalign{}
\begin{minipage}[b]{\linewidth}\raggedright
\#
\end{minipage} & \begin{minipage}[b]{\linewidth}\raggedright
Test
\end{minipage} & \begin{minipage}[b]{\linewidth}\raggedright
Expect
\end{minipage} & \begin{minipage}[b]{\linewidth}\raggedright
Proves
\end{minipage} \\
\midrule\noalign{}
\endhead
\bottomrule\noalign{}
\endlastfoot
2.6.1 & DC volts, each monitored line to GND, all switches
open & 4.8--5.25 V (the rail through 10.47 k) & The board
only ever \emph{pulls up weakly}; it never drives a line \\
2.6.2 & DC \textbf{mA} meter from each monitored line to GND
(meter in place of the switch) & \textbf{0.45--0.55 mA} &
The \textbf{most} current the board can push into a bus line
in normal operation: \textless{} 2 \% of the 31 mA the
terminator already makes the EI sink \\
2.6.3 & DC volts, each \textbf{no-connect} finger to GND
(probe the ribbon conductor) & 0.00 V (floating) & The board
leaves the other 12 signals alone \\
2.6.4 & Board supply current with all five switches closed
(lines held low), no LEDs counted & idle + 5 × 0.48 ≈ idle +
2.4 mA (plus the LEDs that light) & Matches the pull-up
model \\
\end{longtable}

\textbf{2.6.B --- Logic truth table} (jig pull-ups off, then
repeat with them \textbf{on})

\begin{longtable}[]{@{}
  >{\raggedright\arraybackslash}p{(\columnwidth - 6\tabcolsep) * \real{0.2500}}
  >{\raggedright\arraybackslash}p{(\columnwidth - 6\tabcolsep) * \real{0.2500}}
  >{\raggedright\arraybackslash}p{(\columnwidth - 6\tabcolsep) * \real{0.2500}}
  >{\raggedright\arraybackslash}p{(\columnwidth - 6\tabcolsep) * \real{0.2500}}@{}}
\toprule\noalign{}
\begin{minipage}[b]{\linewidth}\raggedright
\#
\end{minipage} & \begin{minipage}[b]{\linewidth}\raggedright
Switches closed
\end{minipage} & \begin{minipage}[b]{\linewidth}\raggedright
Expect lit
\end{minipage} & \begin{minipage}[b]{\linewidth}\raggedright
Must stay dark
\end{minipage} \\
\midrule\noalign{}
\endhead
\bottomrule\noalign{}
\endlastfoot
2.6.5 & none & D13 & D1--D8 \\
2.6.6 & DS1 & D1 & all other drive LEDs \\
2.6.7 & DS2 & D2 & '' \\
2.6.8 & DS3 & D3 & '' \\
2.6.9 & DS4 & D4 & '' \\
2.6.10 & WG only & none & D1--D8 (write gate with no drive
selected lights nothing) \\
2.6.11 & DS1 + WG & D1 + D5 & others \\
2.6.12 & DS2 + WG & D2 + D6 & others \\
2.6.13 & DS3 + WG & D3 + D7 & others \\
2.6.14 & DS4 + WG & D4 + D8 & others \\
2.6.15 & DS1 + DS3 (two selects, not legal on a real bus) &
D1 + D3 & --- shows the channels are independent \\
\end{longtable}

\textbf{2.6.C --- Levels} (one channel, DS1 + WG closed, jig
pull-ups \textbf{on})

\begin{longtable}[]{@{}
  >{\raggedright\arraybackslash}p{(\columnwidth - 4\tabcolsep) * \real{0.3333}}
  >{\raggedright\arraybackslash}p{(\columnwidth - 4\tabcolsep) * \real{0.3333}}
  >{\raggedright\arraybackslash}p{(\columnwidth - 4\tabcolsep) * \real{0.3333}}@{}}
\toprule\noalign{}
\begin{minipage}[b]{\linewidth}\raggedright
\#
\end{minipage} & \begin{minipage}[b]{\linewidth}\raggedright
Measure
\end{minipage} & \begin{minipage}[b]{\linewidth}\raggedright
Expect
\end{minipage} \\
\midrule\noalign{}
\endhead
\bottomrule\noalign{}
\endlastfoot
2.6.16 & J1-10 (finger) & ≈ 0 V (switch) \\
2.6.17 & U1-1 (receiver input, after R17) & \textless{} 0.4
V \\
2.6.18 & U1-2 (DS1, active-high) & \textgreater{} 2.7 V \\
2.6.19 & U3-2 (red cathode) / U2-3 (green cathode) &
\textless{} 0.5 V each \\
2.6.20 & Voltage across R1 (red D1 resistor) & 2.5--2.9 V →
≈ 5.3--6.2 mA \\
2.6.21 & Switch open again: U1-1 & \textgreater{} 4.5 V;
U1-2 \textless{} 0.5 V; LEDs off \\
\end{longtable}

If a channel fails, follow it down that row of table 0.4:
finger → series R → U1 input → U1 output → U3/U2 output →
LED. The first stage that reads wrong is the fault.

\textbf{2.6.D --- Unpowered board on a powered bus} (EI on,
monitor off --- the reverse case)

\begin{longtable}[]{@{}
  >{\raggedright\arraybackslash}p{(\columnwidth - 4\tabcolsep) * \real{0.3333}}
  >{\raggedright\arraybackslash}p{(\columnwidth - 4\tabcolsep) * \real{0.3333}}
  >{\raggedright\arraybackslash}p{(\columnwidth - 4\tabcolsep) * \real{0.3333}}@{}}
\toprule\noalign{}
\begin{minipage}[b]{\linewidth}\raggedright
\#
\end{minipage} & \begin{minipage}[b]{\linewidth}\raggedright
Test
\end{minipage} & \begin{minipage}[b]{\linewidth}\raggedright
Expect
\end{minipage} \\
\midrule\noalign{}
\endhead
\bottomrule\noalign{}
\endlastfoot
2.6.22 & Board supply \textbf{disconnected} from J2. Jig
pull-ups \textbf{on} (+5V\_JIG = 5.0 V). mA meter in series
with one 150 Ω pull-up → line & Line current changes by
\textbf{\textless{} 0.6 mA} between the board plugged in and
the board off the jig \\
2.6.23 & With all five pull-ups on, measure the dead board's
+5V rail (U1-14) & Record. Expect well under 2 V; D13 may
glow very faintly. The back-feed is limited by 470 Ω + 10 kΩ
per line \\
2.6.24 & Line high voltage at the jig with the dead board
plugged in & ≥ 4.9 V (the board doesn't drag the bus down
when it's unpowered) \\
\end{longtable}

\hypertarget{protection-tests-board-1-only}{%
\subsubsection{\texorpdfstring{2.7 Protection tests ---
\textbf{board 1
only}}{2.7 Protection tests --- board 1 only}}\label{protection-tests-board-1-only}}

These deliberately apply the faults in
\texttt{THEORY\_OF\_OPERATION.md} §4.3. \textbf{Pull U1, U2
and U3 from their sockets first} (note the orientation) so
an unexpected result can only cost a cheap part. Jig
disconnected.

\begin{longtable}[]{@{}
  >{\raggedright\arraybackslash}p{(\columnwidth - 8\tabcolsep) * \real{0.2000}}
  >{\raggedright\arraybackslash}p{(\columnwidth - 8\tabcolsep) * \real{0.2000}}
  >{\raggedright\arraybackslash}p{(\columnwidth - 8\tabcolsep) * \real{0.2000}}
  >{\raggedright\arraybackslash}p{(\columnwidth - 8\tabcolsep) * \real{0.2000}}
  >{\raggedright\arraybackslash}p{(\columnwidth - 8\tabcolsep) * \real{0.2000}}@{}}
\toprule\noalign{}
\begin{minipage}[b]{\linewidth}\raggedright
\#
\end{minipage} & \begin{minipage}[b]{\linewidth}\raggedright
Fault applied
\end{minipage} & \begin{minipage}[b]{\linewidth}\raggedright
How
\end{minipage} & \begin{minipage}[b]{\linewidth}\raggedright
Expect
\end{minipage} & \begin{minipage}[b]{\linewidth}\raggedright
Pass
\end{minipage} \\
\midrule\noalign{}
\endhead
\bottomrule\noalign{}
\endlastfoot
2.7.1 & \textbf{Reversed +5 V} & 5.0 V, 50 mA limit,
\textbf{+ on J2-3, − on J2-4} & Current \textless{} 1 mA;
rail (U1 socket pin 14 to GND) ≈ 0 V; D13 dark & Q1
blocks \\
2.7.2 & \textbf{Reversed +12 V} & 12 V, 50 mA limit, + on
J2-2, − on J2-1 & Current \textless{} 0.1 mA; D12 dark & D12
blocks \\
2.7.3 & \textbf{Overvoltage on +5 V} (clamp) & Normal
polarity, \textbf{300 mA limit}. Raise J2-4 from 5.0 V in
0.2 V steps to 7.0 V or until the supply hits its limit. At
each step record the input V, the current and the rail V &
Rail follows the input up to D11's knee (≈ 5.3--5.9 V);
after that the current climbs steeply while the rail rises
slowly & Rail \textbf{\textless{} 6.5 V} at 300 mA (74LS
absolute max is 7.0 V). Don't hold more than 30 s at 300
mA \\
2.7.4 & \emph{Optional} --- \textbf{12 V on the +5 V pin}
(miswired cable) & Normal polarity, \textbf{12 V, 2 A
limit}, ICs still out. Watch the current; switch off after
10 s regardless & Current jumps to ≈ 1--2 A, then F1 trips
and the current falls to tens of mA; rail stays \textless{}
7 V. Let it cool 5 min, then repeat 2.1 → normal & F1 trips;
nothing smokes; afterwards the board passes 2.1--2.2 with
the ICs refitted \\
\end{longtable}

Skip 2.7.4 if you'd rather not stress D11 (5 W part, short
pulse above rating). Steps 2.7.1--2.7.3 already show the
clamp and the reverse block work.

\hypertarget{worst-case-fault-current-into-a-bus-line-board-1-only}{%
\subsubsection{\texorpdfstring{2.8 Worst-case fault current
into a bus line --- \textbf{board 1
only}}{2.8 Worst-case fault current into a bus line --- board 1 only}}\label{worst-case-fault-current-into-a-bus-line-board-1-only}}

This checks the worst case the 470 Ω series resistors are
there for: a dead receiver shorting a bus line to the rail.

\begin{enumerate}
\def\labelenumi{\arabic{enumi}.}
\tightlist
\item
  ICs \textbf{out}. Board at 5.00 V, 100 mA limit.
\item
  Jumper \textbf{U1 socket pin 1 (DS1\_in) to U1 socket pin
  14 (+5 V)}. This simulates a receiver input welded to the
  rail.
\item
  mA meter from \textbf{J1-10} (finger) to \textbf{GND}.
\end{enumerate}

\begin{longtable}[]{@{}
  >{\raggedright\arraybackslash}p{(\columnwidth - 4\tabcolsep) * \real{0.3333}}
  >{\raggedright\arraybackslash}p{(\columnwidth - 4\tabcolsep) * \real{0.3333}}
  >{\raggedright\arraybackslash}p{(\columnwidth - 4\tabcolsep) * \real{0.3333}}@{}}
\toprule\noalign{}
\begin{minipage}[b]{\linewidth}\raggedright
\#
\end{minipage} & \begin{minipage}[b]{\linewidth}\raggedright
Expect
\end{minipage} & \begin{minipage}[b]{\linewidth}\raggedright
Pass
\end{minipage} \\
\midrule\noalign{}
\endhead
\bottomrule\noalign{}
\endlastfoot
2.8.1 & ≈ 10.6 mA (5.0 V / 470 Ω) & \textbf{≤ 12 mA} ---
added to the terminator's 31 mA, that's 43 mA, inside the
7438's 48 mA rating \\
\end{longtable}

Repeat for one more channel (e.g.~U1 pin 11 / J1-24). Remove
the jumper, refit the ICs (check the notches), and re-run
2.1 and 2.6.B quickly.

Phase 2 done → record, go to Phase 3 (board 1 first).

\begin{center}\rule{0.5\linewidth}{0.5pt}\end{center}

\hypertarget{phase-3-on-the-expansion-interface}{%
\subsection{Phase 3 --- On the Expansion
Interface}\label{phase-3-on-the-expansion-interface}}

\hypertarget{before-starting}{%
\subsubsection{3.0 Before starting}\label{before-starting}}

\begin{itemize}
\tightlist
\item[$\square$]
  Board has passed Phase 1 and Phase 2.
\item[$\square$]
  \textbf{Baseline without the monitor:} Model I + EI +
  drives boot DOS, \texttt{DIR} each drive, read a few
  files. Write down anything that already misbehaves, so it
  isn't blamed on the monitor later.
\item[$\square$]
  \textbf{Cable.} Use the six-connector cable from
  \texttt{docs/CABLE.md}, or a straight-through cable with a
  spare connector. \textbf{On a Radio Shack ``pulled-pin''
  cable} each drive connector carries only one drive-select
  line, so the monitor would only ever see that one drive.
  Put it on a connector with all 34 contacts.
\item[$\square$]
  Cable continuity: each connector, pin N to pin N, and no
  adjacent-pin shorts (DMM, 2 minutes).
\item[$\square$]
  Red stripe = pin 1 on every connector. On the monitor, pin
  1 is the solder-side finger directly under pin 2, at the
  \textbf{board-corner end of the tab, away from J2}.
  Confirm against the silk pin numbers before mating.
\item[$\square$]
  The terminated drive is at the \textbf{far end} of the
  cable. The monitor adds no terminator.
\end{itemize}

\hypertarget{ground-and-supply-check}{%
\subsubsection{3.1 Ground and supply
check}\label{ground-and-supply-check}}

\begin{enumerate}
\def\labelenumi{\arabic{enumi}.}
\tightlist
\item
  Model I + EI \textbf{on}, cable connected to the EI but
  \textbf{not} to the monitor. Monitor supply on but not
  connected to the monitor.
\item
  DC volts between the \textbf{bench supply (−)} and the
  \textbf{EI chassis / a GND pin}: \textless{} 0.1 V. AC
  volts: \textless{} 0.1 V.

  \begin{itemize}
  \tightlist
  \item
    A higher reading means the bench supply or the EI is
    floating at a different potential. Connecting it through
    the cable's ground would push current through the
    ribbon. \textbf{Fix first} (bring both to the same mains
    ground, or use a floating supply).
  \end{itemize}
\item
  Everything \textbf{off}.
\end{enumerate}

\hypertarget{stage-a-ei-monitor-no-drives}{%
\subsubsection{3.2 Stage A --- EI + monitor, no
drives}\label{stage-a-ei-monitor-no-drives}}

\begin{enumerate}
\def\labelenumi{\arabic{enumi}.}
\tightlist
\item
  Everything off. Plug the monitor onto the DIAG connector
  (stripe to pin 1). \textbf{No drives connected.}
\item
  Monitor supply: 5.00 V, 100 mA limit (+12 V optional).
  Power on → idle current same as 2.1.1 (± 2 mA).
\item
  Power on the Model I / EI.
\end{enumerate}

\begin{longtable}[]{@{}
  >{\raggedright\arraybackslash}p{(\columnwidth - 4\tabcolsep) * \real{0.3333}}
  >{\raggedright\arraybackslash}p{(\columnwidth - 4\tabcolsep) * \real{0.3333}}
  >{\raggedright\arraybackslash}p{(\columnwidth - 4\tabcolsep) * \real{0.3333}}@{}}
\toprule\noalign{}
\begin{minipage}[b]{\linewidth}\raggedright
\#
\end{minipage} & \begin{minipage}[b]{\linewidth}\raggedright
Observe
\end{minipage} & \begin{minipage}[b]{\linewidth}\raggedright
Expect
\end{minipage} \\
\midrule\noalign{}
\endhead
\bottomrule\noalign{}
\endlastfoot
3.2.1 & Monitor current as the EI powers up & stays within
idle + 15 mA, no spike to the limit \\
3.2.2 & At power-on, with no drive fitted & The ROM tries to
boot drive :0, so \textbf{D1 should light} (DS1 asserted).
Hold BREAK (or reset with BREAK) to get to BASIC \\
3.2.3 & Run the select-cycling program below & D1 → D2 → D3
→ D4 in turn, about 1 s each; \textbf{no green ever} \\
3.2.4 & DMM on J1-10 side of R17 (or scope) while :0
selected & ≤ 0.4 V asserted, ≈ 5 V released \\
\end{longtable}

\begin{Shaded}
\begin{Highlighting}[]
\NormalTok{10 REM CYCLE DRIVE SELECTS :0{-}:3, 37E1H = 14305 (MODEL I DRIVE SELECT LATCH)}
\NormalTok{20 FOR D = 0 TO 3}
\NormalTok{30 V = 2 \^{} D}
\NormalTok{40 FOR T = 1 TO 150 : POKE 14305, V : NEXT T}
\NormalTok{50 NEXT D}
\NormalTok{60 GOTO 20}
\end{Highlighting}
\end{Shaded}

The EI's select latch is cleared by the motor-on timeout
after a few seconds, which is why line 40 keeps re-writing
it. Without drives there's no terminator, so the monitor's
10 k pull-ups are the only thing pulling the lines high.
That's fine for this test.

\hypertarget{stage-b-ei-monitor-drives}{%
\subsubsection{3.3 Stage B --- EI + monitor +
drives}\label{stage-b-ei-monitor-drives}}

\begin{enumerate}
\def\labelenumi{\arabic{enumi}.}
\tightlist
\item
  Everything off. Connect the drives (terminated drive at
  the end). Monitor on DIAG.
\item
  Power order: monitor supply → drives → Model I / EI.
\end{enumerate}

\begin{longtable}[]{@{}
  >{\raggedright\arraybackslash}p{(\columnwidth - 4\tabcolsep) * \real{0.3333}}
  >{\raggedright\arraybackslash}p{(\columnwidth - 4\tabcolsep) * \real{0.3333}}
  >{\raggedright\arraybackslash}p{(\columnwidth - 4\tabcolsep) * \real{0.3333}}@{}}
\toprule\noalign{}
\begin{minipage}[b]{\linewidth}\raggedright
\#
\end{minipage} & \begin{minipage}[b]{\linewidth}\raggedright
Test
\end{minipage} & \begin{minipage}[b]{\linewidth}\raggedright
Expect
\end{minipage} \\
\midrule\noalign{}
\endhead
\bottomrule\noalign{}
\endlastfoot
3.3.1 & Boot DOS from :0 & Boots exactly as in the baseline.
D1 lights whenever drive :0's own activity LED does \\
3.3.2 & \texttt{DIR} on each drive present & Red LED for
that drive only, in step with the drive's own LED \\
3.3.3 & Run the BASIC select-cycling program (from Disk
BASIC) & As 3.2.3; the drive motors start, which is
normal \\
3.3.4 & \textbf{Write test} --- scratch disk in :1. Save a
file or format it (syntax depends on the DOS) & D2 red
throughout, \textbf{D6 green only while writing}. No green
on any other drive \\
3.3.5 & Read back / \texttt{DIR} the scratch disk; copy a
file :0 → :1 and compare & No errors, same as the
baseline \\
3.3.6 & \emph{(scope, optional)} DS1 at the DIAG connector,
monitor plugged in vs removed (everything off to swap) & Low
level differs by \textless{} 50 mV; rise and fall times look
the same \\
\end{longtable}

\hypertarget{power-order-robustness}{%
\subsubsection{3.4 Power-order
robustness}\label{power-order-robustness}}

Run a read loop (e.g.~repeated \texttt{DIR\ :0}, or a BASIC
program that LOADs a file over and over) and,
\textbf{without touching the card edge}:

\begin{longtable}[]{@{}
  >{\raggedright\arraybackslash}p{(\columnwidth - 4\tabcolsep) * \real{0.3333}}
  >{\raggedright\arraybackslash}p{(\columnwidth - 4\tabcolsep) * \real{0.3333}}
  >{\raggedright\arraybackslash}p{(\columnwidth - 4\tabcolsep) * \real{0.3333}}@{}}
\toprule\noalign{}
\begin{minipage}[b]{\linewidth}\raggedright
\#
\end{minipage} & \begin{minipage}[b]{\linewidth}\raggedright
Action
\end{minipage} & \begin{minipage}[b]{\linewidth}\raggedright
Expect
\end{minipage} \\
\midrule\noalign{}
\endhead
\bottomrule\noalign{}
\endlastfoot
3.4.1 & Unplug the monitor's \textbf{J2} (power) while the
loop runs & Loop continues with no errors; monitor LEDs go
dark \\
3.4.2 & Plug J2 back in & Loop continues; LEDs resume \\
3.4.3 & Monitor powered, Model I / EI and drives off &
Monitor current back to idle; no LEDs except D13 (and
D12) \\
3.4.4 & Monitor unpowered, power the system on and boot &
Boots normally (an unpowered monitor doesn't load the
bus) \\
\end{longtable}

\hypertarget{soak}{%
\subsubsection{3.5 Soak}\label{soak}}

\begin{longtable}[]{@{}
  >{\raggedright\arraybackslash}p{(\columnwidth - 4\tabcolsep) * \real{0.3333}}
  >{\raggedright\arraybackslash}p{(\columnwidth - 4\tabcolsep) * \real{0.3333}}
  >{\raggedright\arraybackslash}p{(\columnwidth - 4\tabcolsep) * \real{0.3333}}@{}}
\toprule\noalign{}
\begin{minipage}[b]{\linewidth}\raggedright
\#
\end{minipage} & \begin{minipage}[b]{\linewidth}\raggedright
Test
\end{minipage} & \begin{minipage}[b]{\linewidth}\raggedright
Expect
\end{minipage} \\
\midrule\noalign{}
\endhead
\bottomrule\noalign{}
\endlastfoot
3.5.1 & 1 hour of mixed read/write on a scratch disk & No
disk errors beyond the baseline \\
3.5.2 & Heat check at the end: U1--U3, Q1, D11, F1, R1--R8 &
All barely warm (\textless{} 45 °C). \textbf{D11 warm} means
the rail is above its knee: check the supply voltage \\
3.5.3 & Monitor current at the end & Same as at the start (±
2 mA) \\
\end{longtable}

\hypertarget{final-power-source}{%
\subsubsection{3.6 Final power
source}\label{final-power-source}}

If the monitor will normally be powered from something other
than the bench supply (a PC PSU Molex lead, a drive's
supply), repeat \textbf{2.1--2.2} on that source with a
meter in series (expect 4.75--5.25 V at J2-4), then
\textbf{3.3.1--3.3.4} briefly.

\begin{center}\rule{0.5\linewidth}{0.5pt}\end{center}

\hypertarget{if-something-fails}{%
\subsection{4. If something
fails}\label{if-something-fails}}

\begin{longtable}[]{@{}
  >{\raggedright\arraybackslash}p{(\columnwidth - 4\tabcolsep) * \real{0.3333}}
  >{\raggedright\arraybackslash}p{(\columnwidth - 4\tabcolsep) * \real{0.3333}}
  >{\raggedright\arraybackslash}p{(\columnwidth - 4\tabcolsep) * \real{0.3333}}@{}}
\toprule\noalign{}
\begin{minipage}[b]{\linewidth}\raggedright
Symptom
\end{minipage} & \begin{minipage}[b]{\linewidth}\raggedright
Most likely
\end{minipage} & \begin{minipage}[b]{\linewidth}\raggedright
Look at
\end{minipage} \\
\midrule\noalign{}
\endhead
\bottomrule\noalign{}
\endlastfoot
Current at the limit on first power-up & Reversed IC, solder
bridge, C4 reversed & 1.1.1, 1.1.3, 1.3.4 \\
Rail ≈ 0.5 V below J2-4 & Q1 conducting only through its
body diode & 2.3, Q1 orientation / part \\
All LEDs dead, D13 lit & U1 missing/reversed, or the fingers
not making contact & U1-14 voltage, 2.6.C \\
One red never lights & That row of table 0.4 & Series R,
pull-up, U1 gate, U3 gate, LED polarity \\
Red lights but its green never does & That U2 gate, its LED,
or WG path (R14, U1-11/10) & 2.6.C \\
Green lights with its red even when not writing & WG stuck
active: U1-11 shorted to GND, or U1 gate 5 faulty & U1-11
should be \textgreater{} 4.5 V with WG released; U1-10
\textless{} 0.5 V \\
Green never lights on any drive & WG path open: R14, or
U1-11/U1-10 & 2.6.C on the WG row \\
LEDs flicker with nothing asserted & Pull-up open, so a
floating LS14 input & R9--R13 \\
Disk errors with the monitor fitted but not without & Cable
(the sixth connector), terminator position, or a monitored
line loaded & 3.0 cable checks, 2.6.1--2.6.2 on that board,
3.3.6 \\
\end{longtable}

\begin{center}\rule{0.5\linewidth}{0.5pt}\end{center}

\clearpage\renewcommand{\arraystretch}{2.1}
\hypertarget{record-sheets}{%
\subsection{5. Record sheets}\label{record-sheets}}

Tick ✓ or write the reading. One column per board.

\hypertarget{phase-1}{%
\subsubsection{5.1 Phase 1}\label{phase-1}}

\begin{longtable}[]{@{}
  >{\raggedright\arraybackslash}p{(\columnwidth - 10\tabcolsep) * \real{0.6250}}
  >{\raggedright\arraybackslash}p{(\columnwidth - 10\tabcolsep) * \real{0.0750}}
  >{\raggedright\arraybackslash}p{(\columnwidth - 10\tabcolsep) * \real{0.0750}}
  >{\raggedright\arraybackslash}p{(\columnwidth - 10\tabcolsep) * \real{0.0750}}
  >{\raggedright\arraybackslash}p{(\columnwidth - 10\tabcolsep) * \real{0.0750}}
  >{\raggedright\arraybackslash}p{(\columnwidth - 10\tabcolsep) * \real{0.0750}}@{}}
\toprule\noalign{}
\begin{minipage}[b]{\linewidth}\raggedright
Step
\end{minipage} & \begin{minipage}[b]{\linewidth}\raggedright
B1
\end{minipage} & \begin{minipage}[b]{\linewidth}\raggedright
B2
\end{minipage} & \begin{minipage}[b]{\linewidth}\raggedright
B3
\end{minipage} & \begin{minipage}[b]{\linewidth}\raggedright
B4
\end{minipage} & \begin{minipage}[b]{\linewidth}\raggedright
B5
\end{minipage} \\
\midrule\noalign{}
\endhead
\bottomrule\noalign{}
\endlastfoot
1.1.1--1.1.8 visual & & & & & \\
1.1.9 thickness (mm) & & & & & \\
1.1.10--1.1.11 fit & & & & & \\
1.1.12 tab-corner / key / J2 drill (mm) & & & & & \\
1.2.1 +12↔+5 OL & & & & & \\
1.2.4 J2-4↔GND (Ω) & & & & & \\
1.2.6 F1 (Ω) & & & & & \\
1.3.1 Q1 body diode (V) & & & & & \\
1.3.4 GND↔rail diode (V) & & & & & \\
1.4.1--1.4.2 no-connect fingers OL & & & & & \\
1.4.3 GND fingers & & & & & \\
1.4.4 series R (Ω) 10/12/14/6/24 & & & & & \\
1.4.5 finger→rail (kΩ) 10/12/14/6/24 & & & & & \\
1.4.8 fingers↔+12 OL & & & & & \\
1.5 LEDs & & & & & \\
\end{longtable}

\hypertarget{phase-2}{%
\subsubsection{5.2 Phase 2}\label{phase-2}}

\begin{longtable}[]{@{}
  >{\raggedright\arraybackslash}p{(\columnwidth - 10\tabcolsep) * \real{0.6250}}
  >{\raggedright\arraybackslash}p{(\columnwidth - 10\tabcolsep) * \real{0.0750}}
  >{\raggedright\arraybackslash}p{(\columnwidth - 10\tabcolsep) * \real{0.0750}}
  >{\raggedright\arraybackslash}p{(\columnwidth - 10\tabcolsep) * \real{0.0750}}
  >{\raggedright\arraybackslash}p{(\columnwidth - 10\tabcolsep) * \real{0.0750}}
  >{\raggedright\arraybackslash}p{(\columnwidth - 10\tabcolsep) * \real{0.0750}}@{}}
\toprule\noalign{}
\begin{minipage}[b]{\linewidth}\raggedright
Step
\end{minipage} & \begin{minipage}[b]{\linewidth}\raggedright
B1
\end{minipage} & \begin{minipage}[b]{\linewidth}\raggedright
B2
\end{minipage} & \begin{minipage}[b]{\linewidth}\raggedright
B3
\end{minipage} & \begin{minipage}[b]{\linewidth}\raggedright
B4
\end{minipage} & \begin{minipage}[b]{\linewidth}\raggedright
B5
\end{minipage} \\
\midrule\noalign{}
\endhead
\bottomrule\noalign{}
\endlastfoot
2.1.1 idle current (mA) & & & & & \\
2.2.2 rail (V) & & & & & \\
2.3.1 Q1 drop (V) & & & & & \\
2.4.1 works at 4.75 V & & & & & \\
2.4.2 current at 5.25 V (mA) & & & & & \\
2.5.1 +12 current (mA) & & & & & \\
2.6.1 line open V (min/max) & & & & & \\
2.6.2 line short mA (max) & & & & & \\
2.6.5--2.6.15 truth table, pull-ups off & & & & & \\
2.6.5--2.6.15 truth table, pull-ups on & & & & & \\
2.6.20 LED resistor V & & & & & \\
2.6.22 dead-board line current (mA) & & & & & \\
2.6.23 back-fed rail (V) & & & & & \\
2.7.1 reverse +5 (mA) & & --- & --- & --- & --- \\
2.7.3 rail at 300 mA (V) & & --- & --- & --- & --- \\
2.7.4 PTC trip (optional) & & --- & --- & --- & --- \\
2.8.1 fault current (mA) & & --- & --- & --- & --- \\
\end{longtable}

\hypertarget{phase-3}{%
\subsubsection{5.3 Phase 3}\label{phase-3}}

\begin{longtable}[]{@{}
  >{\raggedright\arraybackslash}p{(\columnwidth - 10\tabcolsep) * \real{0.6250}}
  >{\raggedright\arraybackslash}p{(\columnwidth - 10\tabcolsep) * \real{0.0750}}
  >{\raggedright\arraybackslash}p{(\columnwidth - 10\tabcolsep) * \real{0.0750}}
  >{\raggedright\arraybackslash}p{(\columnwidth - 10\tabcolsep) * \real{0.0750}}
  >{\raggedright\arraybackslash}p{(\columnwidth - 10\tabcolsep) * \real{0.0750}}
  >{\raggedright\arraybackslash}p{(\columnwidth - 10\tabcolsep) * \real{0.0750}}@{}}
\toprule\noalign{}
\begin{minipage}[b]{\linewidth}\raggedright
Step
\end{minipage} & \begin{minipage}[b]{\linewidth}\raggedright
B1
\end{minipage} & \begin{minipage}[b]{\linewidth}\raggedright
B2
\end{minipage} & \begin{minipage}[b]{\linewidth}\raggedright
B3
\end{minipage} & \begin{minipage}[b]{\linewidth}\raggedright
B4
\end{minipage} & \begin{minipage}[b]{\linewidth}\raggedright
B5
\end{minipage} \\
\midrule\noalign{}
\endhead
\bottomrule\noalign{}
\endlastfoot
3.1 ground offset (V DC / AC) & & & & & \\
3.2.1--3.2.4 EI only & & & & & \\
3.3.1--3.3.5 with drives & & & & & \\
3.3.6 scope delta (mV) & & & & & \\
3.4.1--3.4.4 power order & & & & & \\
3.5 soak / heat / current & & & & & \\
3.6 final supply & & & & & \\
\end{longtable}

\begin{center}\rule{0.5\linewidth}{0.5pt}\end{center}

\renewcommand{\arraystretch}{1.25}
\hypertarget{todo}{%
\subsection{6. TODO}\label{todo}}

\begin{itemize}
\tightlist
\item[$\square$]
  Build the bus-simulator jig (0.6)
\item[$\square$]
  Board 1: Phase 1 → Phase 2 (all, including 2.7/2.8) →
  Phase 3
\item[$\square$]
  Boards 2--5: Phase 1 → Phase 2 (2.1--2.6) → Phase 3
  (3.3.1--3.3.4)
\item[$\square$]
  Feed the 1.1.12 measurements into \texttt{gen/config.py}
  and close the ASSUMED items in \texttt{docs/MECHANICAL.md}
\item[$\square$]
  Bring \texttt{docs/BOM.md} and
  \texttt{THEORY\_OF\_OPERATION.md} up to the r2 values
  (table 0.3)
\item[$\square$]
  Results go into the v0.3 plan
  (\texttt{docs/V0.3\_PLAN.md})
\end{itemize}

\end{document}
